%!TEX root = main.tex
\chapter{문법편}

\subsection{일러두기}
\begin{itemize}[noitemsep]
  \item 이미 관찰된 문법 요소는 아래 체크리스트를 확인한 뒤 넘어가십시오.
  \item 예문과 조금 다르더라도 특정 문법 요소가 드러난다면 넘어가도 좋습니다.
  \item 표준형 문법 요소를 억지로 이끌어낼 필요는 없습니다. 다만 표준형 대신 쓰인 문법 요소를 잘 기록하고 넘어가도록 합니다. 이때 방언의 음운적 특징을 잘 반영하여(예: 뭐라카다, 해뿌리고) 그대로 기록합니다.
  \item 특정 문법 요소가 드러나지 않는다면 역질문을 통해 사용 여부를 확인합니다. 만약 사용한다는 답변을 들은 경우, 해당 문법 요소는 어떤 상황에서 어떤 의미로 사용하는지 함께 확인합니다.
\end{itemize}

\subsection{작성 요령}
\begin{itemize}[noitemsep]
  \item 문법 조사표 담당자는 자료제공인 조사와 어휘 조사 과정에서 발화되는 조사에 표시합니다.
  \item 반드시 점화효과를 고려하여 체크합니다. 자료제공인의 답변 직전의 조사자의 질문에 문법 표지가 먼저 등장한 경우, 자료제공인이 해당 문법 표지를 사용하더라도 체크하지 않습니다. <어휘편>에서 조사자의 자유발화 질문에 포함된 문법 표지에 의한 점화 효과는 고려하지 않되, 자유발화 질문이 점화 효과를 일으킨 것이 분명하다고 판단되면, 체크하지 않습니다. 
  \item 문법 표지가 자유발화에서 확인되면 `방언형' 난에 어느 조사 파트에서 어떤 형식으로 확인됐는지 적습니다.
    \begin{itemize}[noitemsep]
      \item 예 1: `벼'를 유도하던 중 `벼를'을 검출한 경우 → `벼' / 벼를
      \item 예 2: 자료제공인 조사 중 `해가지구'를 검출한 경우 → 자료제공인 조사 / 해가지구
    \end{itemize}
  \item 조사 현장에서 최대한 문법 표지를 확인한 뒤, 하루의 조사를 마무리할 때 녹음을 다시 들으며 놓친 문법 표지를 다시 확인합니다.
  \item 조사 과정에서 확인되지 않은 문법 사항은 아래 질문을 참고하여 보충합니다.
\end{itemize}

\newpage
\section{조사}
※ 질문에서 조사 노출을 최소화하기 바랍니다.

\subsection{격조사}
\Gram{
  word={(격조사) -이/가},
  desc={선행 명사의 모음이 상승하는지도 확인(예: /속+이/ → [쇠기])}
}

\Gram{
  word={(격조사) -을/를}
}

\Gram{
  word={(격조사) -에게/게}
}

\Gram{
  word={(격조사) -에}
}

\Gram{
  word={(격조사) -에서}
}

\Gram{
  word={(격조사) -(으)로}
}

\Gram{
  word={(격조사) -와/과}
}

\Gram{
  word={(격조사) -보다}
}

\Gram{
  word={(격조사) -처럼}
}

\Gram{
  word={(격조사) -만큼}
}

\Gram{
  word={(격조사) -더러, -보고}
}

\Gram{
  word={(격조사) -(이)랑}
}

\Gram{
  word={(격조사) -커녕}
}

\Gram{
  word={(격조사) -아/야}
}

\subsection{보조사}
작성중

\subsection{문장 뒤 조사}
작성중

\section{종결어미}
작성중 - 표 매크로 필요

\section{연결어미}
작성중

\section{시제}
작성중

\section{관형형}
작성중

\section{부정}
작성중

\section{사동표현과 피동표현}
작성중

\subsection{사동표현}
작성중

\subsection{피동표현}
작성중

\subsection{보조용언}
